% paper/sections/02c_nondim_curvature.tex
%
% §2 支配方程式の続き（02_governing.tex からの分割）
% 内容：無次元 Navier--Stokes 方程式・界面曲率の定義
%
\subsection{無次元 Navier--Stokes 方程式}
\label{sec:nondim}
% ═══════════════════════════════════════════════════════════════════════════════

本節では，次元を持つ Navier--Stokes 方程式から，その無次元形を段階的に導出する．
このプロセスを理解することは，シミュレーション結果を物理的に解釈する上で不可欠である．

% ─────────────────────────────────────────────────────────────────────────────
\subsubsection{次元系の支配方程式}
% ─────────────────────────────────────────────────────────────────────────────

まず，次元を持つ（アスタリスク $^*$ 付きの）変数を用いて，
非圧縮性流体の運動量保存則（Navier--Stokes 方程式）を記述する．
表面張力項 $\bm{f}_\sigma^*$ の離散化方針は§~\ref{sec:csf}で概説し，
完全な定式化は付録~\ref{app:csf_full}（CSF）で述べる．
\begin{equation}
 \rho^*\!\left(\pd{\bu^*}{t^*}+\bu^*\cdot\bnabla^*\bu^*\right)
 =-\bnabla^* p^* + \bnabla^*\cdot\!\left[\mu^*(\bnabla^*\bu^*+(\bnabla^*\bu^*)^T)\right]
 +\rho^*\bm{g}^* + \bm{f}_\sigma^*
 \label{eq:NS_dim}
\end{equation}
ここで，$\rho^*$ は密度，$\mu^*$ は粘性係数，$\bu^*$ は速度ベクトル，
$p^*$ は圧力，$\bm{f}_\sigma^*$ は表面張力項である．
重力加速度ベクトル $\bm{g}^*$ は，\textbf{鉛直上向きを正}（$\hat{\bm{z}}$ 方向）とする座標系において
\begin{equation}
  \bm{g}^* = -g\hat{\bm{z}}, \qquad g > 0
  \label{eq:gravity_convention}
\end{equation}
と定義する（$g \approx 9.81\,\mathrm{m/s^2}$）．
符号が負であることから，$+\rho^*\bm{g}^*$ の項は鉛直下向きの力 $-\rho^* g\hat{\bm{z}}$ を与える．
この座標規約は本章全体を通して一貫して用いる．

% ─────────────────────────────────────────────────────────────────────────────
\subsubsection{無次元化変数の導入}
% ─────────────────────────────────────────────────────────────────────────────

次に，参照量として液相密度 $\rho_l$，代表速度 $U$，代表長さ $L$ を選ぶ~\cite{Tryggvason2011}．
$L$ は液滴半径・気泡初期直径・毛管長のいずれかを問題設定に応じて採用し，
$U$ は初速または重力終端速度を採用する（§~\ref{sec:validation} の各ベンチマーク参照）．
本稿の表記規約として，\textbf{*-付きの変数は有次元量，*-無しは無次元量}とする．
これらを用いて，物理変数を以下のように無次元化する．
\begin{align}
 \bm{x}^* &= L\bm{x} & \bu^* &= U\bu & t^* &= \frac{L}{U}t & p^* &= \rho_l U^2 p \\
 \rho^* &= \rho_l \rho(\psi) & \mu^* &= \mu_l \mu(\psi) & \bnabla^* &= \frac{1}{L}\bnabla \\
 \phi^* &= L\phi & \kappa^* &= \kappa/L
\end{align}
ここで，各式の左辺（*-付き）は有次元，右辺（*-無し）は無次元である．
$\phi$ は長さ次元を持つため $\phi^* = L\phi$，$\kappa$ は長さの逆次元を持つため
$\kappa^* = \kappa/L$ と変換され，式~\eqref{eq:curvature_div} の無次元化と整合する．

% ─────────────────────────────────────────────────────────────────────────────
\subsubsection{なぜ「慣性力スケール」で正規化するのか}
% ─────────────────────────────────────────────────────────────────────────────

無次元化の本質的な目的は，「どの力が支配的か」を一目で判断できる形に
方程式を書き換えることである．

式~\eqref{eq:NS_dim} の各項は次元をもつため，
項の大きさを比較しても値が使用する単位系に依存してしまい，
普遍的な議論ができない．
そこで，全ての項を同じ「力のスケール」で割ることで，
各項の係数が純粋な力の比（すなわち無次元数）として現れるようにする．

慣性力スケール $\rho_l U^2/L$ を基準として方程式の両辺を割ることで，
左辺（慣性項）の係数がちょうど $1$ になり，右辺の各項は「慣性力に対する相対的な強さ」として
$Re, Fr, We$ の形の無次元係数に変換される．
慣性力は運動量方程式の基本項であるため，これを基準とするのが標準的な慣行である．

% ─────────────────────────────────────────────────────────────────────────────
\subsubsection{無次元方程式の導出}
% ─────────────────────────────────────────────────────────────────────────────

上記の無次元化変数を式~\eqref{eq:NS_dim} の各項に代入すると，
全ての項に共通因子 $[\rho_l U^2/L]$（慣性力スケール）が現れる．
粘性項からは $1/Re$，重力項からは $1/Fr^2$，表面張力項からは $1/We$ が
自然に析出される（各項の変換過程は付録~\ref{app:nondim_details} 参照）．
両辺を共通因子で割ると，最終的な無次元 Navier--Stokes 方程式が得られる．

\begin{equation}
 \rho(\psi)\!\left(\pd{\bu}{t}+\bu\cdot\bnabla\bu\right)
 =-\bnabla p+\frac{1}{Re}\,\bnabla\cdot\!\left[\mu(\psi)\,(\bnabla\bu+\bnabla\bu^T)\right]
 -\frac{\rho(\psi)}{Fr^2}\hat{\bm z}+\frac{\kappa\,\bnabla\psi}{We}
 \label{eq:NS_dimless_cls}
\end{equation}

この方程式は，流体の挙動が $Re, Fr, We$ という3つの無次元パラメータによって
支配されることを明示している．

各無次元数の定義と物理的意味を表~\ref{tab:dimless_numbers} に示す．

\begin{table}[htbp]
\centering
\renewcommand{\arraystretch}{\PaperTableArrayStretch}
\begin{tabular}{lccc}
\hline
無次元数      & 定義 & 力の比 & 大きいとき \\
\hline
Reynolds 数 $Re$ & $\dfrac{\rho_l U L}{\mu_l}$
  & $\dfrac{\text{慣性力}}{\text{粘性力}}$
  & 乱流的，粘性の影響が小さい \\[6pt]
Froude 数 $Fr$   & $\dfrac{U}{\sqrt{gL}}$
  & $\sqrt{\dfrac{\text{慣性力}}{\text{重力}}}$
  & 重力波が生じにくい（速い流れ） \\[6pt]
Weber 数 $We$    & $\dfrac{\rho_l U^2 L}{\sigma}$
  & $\dfrac{\text{慣性力}}{\text{表面張力}}$
  & 液滴が変形・分裂しやすい \\[6pt]
Capillary 数 $Ca$ & $\dfrac{\mu_l U}{\sigma} = \dfrac{We}{Re}$
  & $\dfrac{\text{粘性力}}{\text{表面張力}}$
  & 粘性流で界面が引き伸ばされる \\[6pt]
Bond 数 $Bo$     & $\dfrac{\rho_l g L^2}{\sigma} = \dfrac{We}{Fr^2}$
  & $\dfrac{\text{重力}}{\text{表面張力}}$
  & 重力変形が表面張力保持を上回る \\
\hline
\end{tabular}
\caption{無次元数の定義と力の比としての解釈}
\label{tab:dimless_numbers}
\PaperCaptionNote{$Ca, Bo$ は独立パラメータではなく $We/Re, We/Fr^2$ から派生する．}
\end{table}

\noindent\textbf{例：}
水（$\rho_l=1000\,\mathrm{kg/m^3}$,
$\mu_l=10^{-3}\,\mathrm{Pa\cdot s}$,
$\sigma=0.073\,\mathrm{N/m}$）の落下液滴
（$U=1\,\mathrm{m/s}$, $L=10^{-3}\,\mathrm{m}$）では，
$Re=1000$，$Fr=10.1$，$We=13.7$，
$Ca = We/Re \approx 1.37\times 10^{-2}$，
$Bo = We/Fr^2 \approx 0.134$ となる．
$We > 1$ であるから，表面張力よりも慣性力が勝り，
液滴は球形から変形しやすい状態にある．
一方 $Ca \ll 1$ であるため粘性応力に対しては表面張力が支配的で，
寄生流れの相対強度は$\Ord{Ca}$ に比例することが知られている~\cite{Popinet2009}．
$Bo < 1$ は重力が球形を強く崩すほどではないことを示すが，
マイクロ流体や微小液滴ではさらに小さな $Bo$ が典型的である．

% ═══════════════════════════════════════════════════════════════════════════════
\subsection{界面曲率の定義と解析表現}
\label{sec:curvature}
% ═══════════════════════════════════════════════════════════════════════════════

無次元 NS 方程式~\eqref{eq:NS_dimless_cls} の表面張力項 $\kappa\bnabla\psi/We$ に現れる
界面の曲率 $\kappa$ は，表面張力項を評価するために不可欠な幾何学量である．
その定義は，外向き法線ベクトル $\hat{\bm{n}}_{lg} = \bnabla\phi/|\bnabla\phi|$（§~\ref{sec:notation}，本稿全章共通）の発散として与えられる~\cite{OsherSethian1988}．
ここではまず $\phi$ に基づく表現を導出し，
後に $\psi$ からの直接計算が可能であることを示す（§~\ref{sec:curvature_direct_psi}）．
\begin{equation}
 \kappa_{lg}(\phi) = \bnabla\cdot \hat{\bm{n}}_{lg}
 = \bnabla\cdot\!\left(\frac{\bnabla\phi}{|\bnabla\phi|}\right)
 \label{eq:curvature_div}
\end{equation}
$\phi$ が真の符号付き距離関数であれば $|\bnabla\phi|=1$ であるため，
$\kappa_{lg} = \bnabla^2\phi$ と単純化される．
しかし，数値計算では移流誤差によりこの性質が常に満たされるとは限らないため，
式~\eqref{eq:curvature_div}を直接扱う必要がある．

% ─────────────────────────────────────────────────────────────────────────────
\subsubsection{2 次元デカルト座標における曲率式}
% ─────────────────────────────────────────────────────────────────────────────

式~\eqref{eq:curvature_div} を2次元 $(x,y)$ 系で展開すると（商の微分法則による導出は
付録~\ref{app:curvature_2d_derivation}），以下の実装に適した形式が得られる：
\begin{equation}
  \kappa_{lg} = \frac{
    \phi_y^2\,\phi_{xx}
    - 2\phi_x\phi_y\,\phi_{xy}
    + \phi_x^2\,\phi_{yy}
  }{\bigl(\phi_x^2 + \phi_y^2\bigr)^{3/2}}
  \label{eq:curvature_2d}
\end{equation}
各微分項（$\phi_x, \phi_y, \phi_{xx}, \phi_{yy}, \phi_{xy}$）を
CCD 法（第~\ref{sec:CCD}章）で評価することで，滑らかな領域では高次精度の曲率が得られる~\cite{MacklinLowengrub2006}．
半径 $R$ の液滴円（$\phi = r - R$）では $\kappa_{lg} = 1/R$ が解析的に再現される
（§~\ref{sec:verify_curvature} で数値的に確認）．

% ─────────────────────────────────────────────────────────────────────────────
\subsubsection{平滑化ヘヴィサイド関数 \texorpdfstring{$\psi$}{ψ} による曲率の直接計算}
\label{sec:curvature_direct_psi}
% ─────────────────────────────────────────────────────────────────────────────

前節で導出した曲率式~\eqref{eq:curvature_2d} は
符号付き距離関数 $\phi$ の偏微分で表されている．
CLS 法では $\psi$ が基本変数であるため，
従来は毎ステップ ロジット逆変換（§~\ref{sec:newton_inversion}）で
$\phi$ を復元してから曲率を計算していた．
しかし，以下に示す\textbf{単調変換に対する曲率の不変性}により，
$\phi$ への逆変換を経ずに $\psi$ の微分値から直接曲率を計算できる．

\begin{theorem}[レベルセット曲率の単調変換不変性]
\label{thm:curvature_invariance}
$I \subset \mathbb{R}$ を開区間とし，$\psi = g(\phi)$，
$g : I \to \mathbb{R}$ は $C^2$ 級の狭義単調減少関数（$g' < 0$ on $I$）とする．
このとき $\phi(\bm{x}) \in I$ かつ $|\bnabla\phi(\bm{x})| \neq 0$ を満たす点 $\bm{x}$ において，
\begin{equation}
  -\bnabla \cdot \frac{\bnabla \psi}{|\bnabla \psi|}
  = \bnabla \cdot \frac{\bnabla \phi}{|\bnabla \phi|}
  = \kappa_{lg}
\end{equation}
すなわち，本稿の $\psi=H_\varepsilon(-\phi)$ 規約では
$-\bnabla\psi/|\bnabla\psi|$ が液相$\to$気相法線を与え，
$\psi$ と $\phi$ のいずれから計算しても同一の向き付き曲率を与える．
\end{theorem}
\begin{proof}
連鎖律より $\bnabla \psi = g'(\phi)\,\bnabla \phi$ である．
$\phi \in I$ において $g' < 0$ より $|g'(\phi)| = -g'(\phi)$ であるから，
\[
  |\bnabla \psi| = -g'(\phi)\,|\bnabla \phi|
\]
したがって，単位法線ベクトルは
\[
  \hat{\bm{n}}_\psi
  \equiv -\frac{\bnabla \psi}{|\bnabla \psi|}
  = -\frac{g'(\phi)\,\bnabla \phi}{-g'(\phi)\,|\bnabla \phi|}
  = \frac{\bnabla \phi}{|\bnabla \phi|}
  \equiv \hat{\bm{n}}_{lg}
\]
となり，$\hat{\bm{n}}_\psi = \hat{\bm{n}}_{lg}$ が成立する．
曲率は $\kappa_{lg} = \bnabla \cdot \hat{\bm{n}}_{lg}$ であるから，
$\psi$ から計算しても $\phi$ から計算しても同一の値を与える．
\end{proof}

\begin{corollary}[CLS ヘヴィサイド関数への適用域]
\label{cor:cls_curvature_domain}
CLS 法の平滑化ヘヴィサイド関数
$g(\phi) = H_\varepsilon(-\phi)=1/(1+e^{\phi/\varepsilon})$ は
$\phi \in \mathbb{R}$ 全域で $C^2$ かつ $g' < 0$ を満たすため，
連続レベルでは定理~\ref{thm:curvature_invariance} が全域で成立する．
一方，$|g'(\phi)| = \varepsilon^{-1}\psi(1-\psi)$ であるから，
$\psi \to 0$ または $\psi \to 1$ の\textbf{飽和領域}では $g'(\phi) \to 0$ に漸近し，
数値離散では $|\bnabla\psi|$ がアンダーフローして曲率の分母が不安定化する．
したがって数値実装では適用域を
\[
  I_\psi \equiv \{\phi : \psi_{\min} < H_\varepsilon(-\phi) < 1 - \psi_{\min}\}
  \quad (\text{推奨 } \psi_{\min} = 0.01)
\]
に限定する．この $\psi_{\min}=0.01$ はロジット写像
$\phi = \varepsilon\,\ln\bigl((1-\psi)/\psi\bigr)$ のもとで
$|\phi| \lesssim 4.6\,\varepsilon$ に対応し（$\ln(0.99/0.01)\approx 4.595$），
CLS の推奨界面幅 $\phi\in[-3\varepsilon,3\varepsilon]$ を包含する最小裕度である．
$H_\varepsilon'(\phi) = \varepsilon^{-1}\psi(1-\psi)$ の飽和裾を最小限除外しつつ
界面近傍帯を全て含む実用的閾値として CLS 文献で慣用される．
具体的な安定化処置は §~\ref{sec:curvature_direct_psi_cls}（第~\ref{sec:levelset}章）で述べる．
\end{corollary}

\noindent
$\phi$ と $\psi$ は系~\ref{cor:cls_curvature_domain} の
$I_\psi$ 上で数学的に等価な曲率を与え，
その範囲では $\psi$（平滑ヘヴィサイド関数）から直接計算する方が有利である：
$\psi$ は移流後も $[0,1]$ に値が収まり再初期化による Eikonal 条件 $|\bnabla\phi|=1$
への依存が弱いため，界面近傍で滑らかな単調プロファイルが保たれている限り数値安定性が高い
（詳細は §~\ref{sec:curvature_direct_psi_cls}）．

\paragraph{曲率の数値的安定化（実装クイックスタート）}
式~\eqref{eq:curvature_2d} の分母 $(\phi_x^2+\phi_y^2)^{3/2}$ は，
移流・再初期化の数値誤差により $|\bnabla\phi|\to 0$ となる格子点でゼロ除算を引き起こす恐れがある．
実装では必ず以下の\textbf{分母クランプ}（下限値 $\varepsilon_{\text{norm}}$）を適用する：
\begin{equation}
  \kappa_{lg} \approx \frac{
    \phi_y^2\,\phi_{xx} - 2\phi_x\phi_y\,\phi_{xy} + \phi_x^2\,\phi_{yy}
  }{\max\!\bigl(\phi_x^2 + \phi_y^2,\; \varepsilon_{\text{norm}}^2\bigr)^{3/2}}
  \label{eq:curvature_clamped}
\end{equation}
（式~\eqref{eq:curvature_2d} と同符号；円 $\phi=r-R$ で $\kappa_{lg}=1/R>0$ を保つ）．
\textbf{推奨値：} $\varepsilon_{\text{norm}} = 10^{-3}$（次元なし，再初期化が正常なら作動しない）．
CLS 再初期化（第~\ref{sec:levelset}章）が $|\bnabla\phi|\approx 1$ を維持していれば，
分母は常に $\approx 1$ に保たれクランプは実質的に不要となる．
クランプは過渡的な状況への防御的フォールバックである．
詳細分析（ゼロ除算の発生メカニズム，誤差評価，ガウシアンフィルタ，実装チェックリスト）は
§~\ref{sec:curvature_stab} を参照．
